\section{Density-aware batch scheduling}
\label{app:density}

The main text (\Cref{sec:method}) reports the wall-clock efficiency of the
cross-document training pipeline as a single measured number. This appendix
documents the mechanism behind it: a data-scheduling scheme that equalizes
attention cost \emph{across data-parallel ranks at every optimizer step}, so the
block-sparse backward pass does not stall the gradient all-reduce. It is a
data-parallel analogue of the intra-sequence load balancing done by
sequence-parallel kernels such as Striped Attention \citep{brandon2023striped}
and DistFlashAttn \citep{li2023distflashattn}: those rebalance the causal
triangle \emph{within} one sequence spread over many devices, whereas we
rebalance the cost of \emph{whole packed sequences} across the ranks of a
data-parallel group. The bucketing borrows the machinery of cost- and
difficulty-ordered curricula
\citep{pouransari2024datasetdecomposition,bengio2009curriculum,kocmi2017curriculum},
but the ordering objective is compute balance, not a learning schedule.

\subsection{The problem: mask sparsity varies across packs}

Under cross-document linking each training example is one packed sequence of $T$
tokens whose attention mask is block-sparse: a query block attends to a key
block only if the two lie in the same document or are joined by a granted link
(\Cref{app:kernels}). The number of live blocks is not constant --- a pack of
many short unlinked documents is nearly block-diagonal, while a pack dominated by
a few long, cross-linked documents is much denser --- and spans roughly a
$6\times$ range across the packed corpus. Because the backward pass runs only
over live blocks, its cost is approximately proportional to that count. Under DDP
with gradient accumulation, every rank must reach the all-reduce together, so the
step is gated by the densest pack in the group; ranks holding lighter packs idle
at the collective. This is the straggler pathology sequence-parallel schemes
solve inside a sequence, resurfacing at the granularity of the data-parallel
batch.

\subsection{The \texttt{kv\_block\_count} metric and its analytic proxy}

Pack density is quantified by \texttt{kv\_block\_count}: the number of non-empty
$(q_\text{block}, k_\text{block})$ pairs in the block-sparse mask on a $128$-token
block grid, a direct proxy for backward cost. It admits two estimators
(\Cref{tab:density-methods}). \emph{Method~B} materializes the actual
\texttt{BlockMask} on the GPU and sums its partial- and full-block counts; this
is the ground truth the kernel sees but runs sequentially, and is used only for
verification behind a GPU-pass flag. \emph{Method~C}, the default, computes the
live-block set analytically from link positions on CPU inside the parallel epoch
precompute workers. The two per-pack costs quoted for the estimators are
$\sim 36$\,ms (Method~B) versus $\sim 1$\,ms (Method~C, at $T=32$k), a
$\sim 36\times$ reduction; these are the estimators' stated costs, not benchmarked
here on fixed hardware.

Method~C accumulates a set of live block pairs in $O(\#\text{blocks})$: for each
document span it adds the lower-triangular block pairs within the span's block
range (intra-document causal), and for each detected link it adds the rectangle
of block pairs formed by the granting query range --- from the link position, or
from the source-span start under a whole-document grant --- against the target
span's block range (cross-document grants). It returns the set's cardinality; the
kernel-level meaning of these block pairs is given in \Cref{app:kernels}.

\paragraph{Conditional exactness.}
Method~C is \emph{exact} with respect to Method~B whenever every grant target ends
before its link, i.e.\ all granted links point backward in packed order. This is
exactly what the packer produces: the target-first packing order
(\Cref{app:traversal}) topologically sorts each connected component so linked-to
documents precede their linkers and emits every component contiguously. Under
that order the analytic and materialized counts agree, and the GPU-pass flag
confirms the equality empirically. The argument covers only acyclic grants; for
cyclic link structures the insertion-order fallback can order some edges forward
and the two counts may diverge, an effect we do not quantify.

\begin{table}[t]
\centering
\small
\begin{tabular}{lllll}
\toprule
Method & Source & Cost/pack & Parallelism & Role \\
\midrule
B & Materialized \texttt{BlockMask} (GPU) & $\sim 36$\,ms & sequential & verification \\
C & Analytic, from link positions (CPU) & $\sim 1$\,ms & per-worker parallel & \textbf{default} \\
\bottomrule
\end{tabular}
\caption{Two estimators for \texttt{kv\_block\_count}. Method~C is the default and
is exact relative to Method~B when every target ends before its link, which the
target-first packing order guarantees. The per-pack costs are the estimators'
stated figures, not benchmarked here.}
\label{tab:density-methods}
\end{table}

\subsection{Quantile bucketing and the bucket-shuffle schedule}

Given a \texttt{kv\_block\_count} for every pack in an epoch, the packs are sorted
by density and the $i$-th of $N$ is assigned bucket
$\lfloor i \cdot n_\text{buckets} / N \rfloor$. These are \emph{equal-count}
quantile buckets, not equal-width density bins, so each bucket holds the same
number of packs regardless of the density distribution's shape; the default is
$n_\text{buckets}=32$. Ties are split by sort position. When corpora are merged,
buckets are re-assigned over the union, since bucket $B$ in one source is not the
same density as bucket $B$ in another.

Buckets are visited in an order balancing two goals: each bucket should be
visited equally often (no density regime over-sampled), yet consecutive steps
should see \emph{different} densities (gradient not correlated with density over
short windows). The visit sequence is built by seeding an RNG with the epoch
index, shuffling the list of buckets, and concatenating $1000$ such shuffles.
Over the full sequence each bucket appears exactly once per shuffle (balance);
within any single shuffle no bucket repeats, so adjacent steps almost always draw
distinct buckets (diversity). Seeding by epoch index alone makes the schedule
deterministic and identical across ranks and restarts.

\subsection{The key invariant: per-step density match across ranks}

The central contribution is what happens at each accumulation step. At global
accumulation step $s$ every rank selects the \emph{same} bucket
\[
  B \;=\; \texttt{bucket\_seq}\big[\, s \bmod \lvert \texttt{bucket\_seq}\rvert \,\big],
\]
because \texttt{bucket\_seq} is seeded by the epoch index and \emph{not} by rank.
Within that bucket, rank $r$ takes the pack at index
$\texttt{consumed}[B] + r$, and the counter is then advanced by
\texttt{world\_size}. Thus at every step all ranks draw packs from one density
bucket, and --- because the bucket is equal-count and density is nearly constant
within a quantile --- their per-rank backward costs are matched. The straggler
gap at the all-reduce closes with no run-time cost measurement at all.

A consequence is that \texttt{world\_size} is \emph{not} baked into the
precompute: the metric, the buckets, and the shuffle sequence are all
rank-count-independent, and only the per-step index arithmetic
($\texttt{consumed}+r$, then $\texttt{consumed}\mathrel{+}=\texttt{world\_size}$)
references it. A run therefore resumes at a different world size from a saved
schedule position --- epoch index, global accumulation step, and per-bucket
consumed counts --- without recomputing anything, the same portability exploited
for the optimizer state. When a bucket lacks \texttt{world\_size} remaining packs
mid-epoch, the schedule scans outward to the nearest bucket that does, in
deterministic order of $\lvert b-B\rvert$; the epoch tail is handled with
\texttt{drop\_last}.

\subsection{Live versus precomputed paths}

Density-aware scheduling is available only on the precomputed data path. The
\emph{live} path builds every pack's block mask on the GPU at each step: it
traverses the document graph, detects links, and constructs the mask online. The
\emph{precomputed} path does link detection once, offline, in parallel workers,
caches the link positions and target ids per pack, and reattaches them at load
time as the baked-grant input to the mask builder (\Cref{app:kernels}). Because
detection, traversal, and density scheduling have all moved offline, the loader
runs synchronously with sub-millisecond load latency; the online cost eliminated
is exactly the per-step link detection, mask build, and graph traversal.

\subsection{Caveats and threats to validity}

None of the following affects correctness --- the mask the kernel executes is
unchanged --- only the tightness of the cost balance, which is consistent with
the measured speedup being close to, but below, the theoretical ceiling in the
main text.

\begin{itemize}
  \item \textbf{Block-size mismatch.} The proxy and the \texttt{BlockMask} count
  blocks on a $128$-token grid, but the production Triton backward kernel runs at
  a $64$-token block size. The metric is therefore \emph{rank-preserving} (denser
  packs stay denser) rather than a literal count of the blocks the executed
  kernel touches.
  \item \textbf{\texttt{max\_grants}.} The analytic count ignores the runtime
  \texttt{max\_grants} cap (default $64$), which drops the lowest-priority links
  at execution time. For packs with more grants than the cap the proxy
  \emph{overcounts} density relative to what the kernel runs; during
  \texttt{max\_grants} warmup the frozen bucket densities also reflect the final
  cap, not the smaller cap in force early in training.
  \item \textbf{Epoch-tail degradation.} The equal-density guarantee holds in the
  epoch interior. Once a bucket can no longer supply \texttt{world\_size} packs,
  the outward scan draws from a neighbouring, differently dense bucket, and
  \texttt{drop\_last} discards up to $\texttt{world\_size}-1$ packs per bucket, so
  balance loosens over the last few steps.
\end{itemize}
